\documentclass[11pt,openany]{book}

% =========================================================
% Packages
% =========================================================
\usepackage[margin=1in]{geometry}

\usepackage{amsmath}
\usepackage{amssymb}
\usepackage{bm}
\usepackage{physics}

\usepackage{graphicx}
\usepackage{booktabs}
\usepackage{enumitem}

\usepackage[
    colorlinks=true,
    linkcolor=blue,
    citecolor=blue,
    urlcolor=blue
]{hyperref}

% =========================================================
% Title Information
% =========================================================
\title{High-Speed I/O Design Notes}
\author{Hanchen Kang}
\date{}

% =========================================================
% Document
% =========================================================
\begin{document}

% =========================================================
% Front Matter
% =========================================================
\frontmatter

\maketitle

\tableofcontents

% =========================================================
% Main Matter
% =========================================================
\mainmatter

% =========================================================
% Chapter 1
% =========================================================
\chapter{Signal Integrity}

\section{Overview}

High-speed I/O channels introduce attenuation, reflection,
dispersion, and inter-symbol interference (ISI).

A transmitted symbol is therefore distorted by the channel
before it reaches the receiver.

The channel can be modeled as a linear time-invariant system
with impulse response

\begin{equation}
h_{\mathrm{ch}}(t).
\end{equation}

If the transmitted pulse is

\begin{equation}
x_{\mathrm{pulse}}(t),
\end{equation}

then the channel pulse response is

\begin{equation}
p_{\mathrm{ch}}(t)
=
x_{\mathrm{pulse}}(t)
*
h_{\mathrm{ch}}(t).
\end{equation}

For a bandwidth-limited channel, the pulse response usually
contains a large main cursor and a long post-cursor tail.

% =========================================================
% Chapter 2
% =========================================================
\chapter{Equalization}

\section{Overview}

Equalization is used to compensate for channel loss and reduce
inter-symbol interference.

Common equalization techniques include:

\begin{itemize}
    \item Continuous-Time Linear Equalizer (CTLE)
    \item Feed-Forward Equalizer (FFE)
    \item Decision-Feedback Equalizer (DFE)
\end{itemize}

The CTLE operates in continuous time and compensates for
frequency-dependent channel attenuation.

The DFE operates using previously detected symbols and is
primarily used to cancel post-cursor ISI.

% =========================================================
% CTLE
% =========================================================
\section{Continuous-Time Linear Equalizer}

\subsection{CTLE Transfer Function}

Consider a typical first-order CTLE transfer function

\begin{equation}
H_{\mathrm{CTLE}}(s)
=
A_0
\frac{1+s/\omega_z}
     {1+s/\omega_p},
\end{equation}

where

\begin{equation}
\omega_z < \omega_p.
\end{equation}

The zero appears at a lower frequency than the pole, producing
high-frequency peaking.

At low frequency,

\begin{equation}
\omega \ll \omega_z,
\end{equation}

the gain is approximately

\begin{equation}
H(j\omega)
\approx
A_0.
\end{equation}

At high frequency,

\begin{equation}
\omega \gg \omega_p,
\end{equation}

the magnitude approaches

\begin{equation}
\left|
H(j\omega)
\right|
\approx
A_0
\frac{\omega_p}{\omega_z}.
\end{equation}

Therefore, the high-frequency boost relative to the low-frequency gain is

\begin{equation}
G_{\mathrm{boost}}
=
20\log_{10}
\left(
\frac{\omega_p}{\omega_z}
\right).
\end{equation}

Thus,

\begin{equation}
\boxed{
\text{CTLE compensates channel high-frequency attenuation.}
}
\end{equation}

% =========================================================
% CTLE Impulse Response
% =========================================================
\subsection{CTLE Impulse Response}

Rewrite the transfer function as

\begin{equation}
H_{\mathrm{CTLE}}(s)
=
A_0
\frac{\omega_p}{\omega_z}
\frac{s+\omega_z}
     {s+\omega_p}.
\end{equation}

Using

\begin{equation}
\frac{s+\omega_z}
{s+\omega_p}
=
1+
\frac{\omega_z-\omega_p}
{s+\omega_p},
\end{equation}

we obtain

\begin{equation}
H_{\mathrm{CTLE}}(s)
=
A_0
\frac{\omega_p}{\omega_z}
\left[
1+
\frac{\omega_z-\omega_p}
{s+\omega_p}
\right].
\end{equation}

Taking the inverse Laplace transform,

\begin{equation}
h_{\mathrm{CTLE}}(t)
=
A_0
\frac{\omega_p}{\omega_z}
\delta(t)
+
A_0
\frac{\omega_p}{\omega_z}
(\omega_z-\omega_p)
e^{-\omega_p t}
u(t).
\end{equation}

Since

\begin{equation}
\omega_p>\omega_z,
\end{equation}

the exponential term has a negative coefficient.

Define

\begin{equation}
K
=
A_0
\frac{\omega_p}{\omega_z},
\end{equation}

and

\begin{equation}
C
=
A_0
\frac{\omega_p}{\omega_z}
(\omega_p-\omega_z),
\end{equation}

where

\begin{equation}
K>0,
\qquad
C>0.
\end{equation}

The CTLE impulse response therefore becomes

\begin{equation}
\boxed{
h_{\mathrm{CTLE}}(t)
=
K\delta(t)
-
Ce^{-\omega_p t}
u(t)
}
\end{equation}

This consists of

\begin{equation}
\boxed{
\text{positive instantaneous response}
+
\text{negative exponentially decaying tail}.
}
\end{equation}

% =========================================================
% Channel Pulse Response
% =========================================================
\subsection{Channel Pulse Response}

Let the transmitted pulse be

\begin{equation}
x_{\mathrm{pulse}}(t).
\end{equation}

The channel pulse response is

\begin{equation}
p_{\mathrm{ch}}(t)
=
x_{\mathrm{pulse}}(t)
*
h_{\mathrm{ch}}(t).
\end{equation}

Because the channel is bandwidth limited, the received pulse
is generally broadened in time.

The response therefore contains

\begin{equation}
\boxed{
\text{main cursor}
+
\text{pre-cursor ISI}
+
\text{post-cursor ISI}.
}
\end{equation}

% =========================================================
% CTLE + Channel
% =========================================================
\subsection{Pulse Response After the CTLE}

After the CTLE, the new pulse response is

\begin{equation}
p_{\mathrm{eq}}(t)
=
p_{\mathrm{ch}}(t)
*
h_{\mathrm{CTLE}}(t).
\end{equation}

Using the continuous-time convolution definition,

\begin{equation}
p_{\mathrm{eq}}(t)
=
\int_{-\infty}^{+\infty}
p_{\mathrm{ch}}(\tau)
h_{\mathrm{CTLE}}(t-\tau)
d\tau.
\end{equation}

Substituting

\begin{equation}
h_{\mathrm{CTLE}}(t)
=
K\delta(t)
-
Ce^{-\omega_p t}
u(t),
\end{equation}

gives

\begin{align}
p_{\mathrm{eq}}(t)
&=
\int_{-\infty}^{+\infty}
p_{\mathrm{ch}}(\tau)
\left[
K\delta(t-\tau)
-
Ce^{-\omega_p(t-\tau)}
u(t-\tau)
\right]
d\tau.
\end{align}

Separate the two terms:

\begin{align}
p_{\mathrm{eq}}(t)
&=
K
\int_{-\infty}^{+\infty}
p_{\mathrm{ch}}(\tau)
\delta(t-\tau)
d\tau
\\
&\quad
-
C
\int_{-\infty}^{+\infty}
p_{\mathrm{ch}}(\tau)
e^{-\omega_p(t-\tau)}
u(t-\tau)
d\tau.
\end{align}

Using the sifting property,

\begin{equation}
\int_{-\infty}^{+\infty}
p_{\mathrm{ch}}(\tau)
\delta(t-\tau)
d\tau
=
p_{\mathrm{ch}}(t).
\end{equation}

Also,

\begin{equation}
u(t-\tau)=0
\qquad
\text{for}
\qquad
\tau>t.
\end{equation}

Therefore,

\begin{equation}
\boxed{
p_{\mathrm{eq}}(t)
=
Kp_{\mathrm{ch}}(t)
-
C
\int_{-\infty}^{t}
p_{\mathrm{ch}}(\tau)
e^{-\omega_p(t-\tau)}
d\tau
}
\end{equation}

This is the most important continuous-time expression for
understanding CTLE equalization.

% =========================================================
% Interpretation
% =========================================================
\subsection{Interpretation of the Integral Term}

Define

\begin{equation}
m(t)
=
\int_{-\infty}^{t}
p_{\mathrm{ch}}(\tau)
e^{-\omega_p(t-\tau)}
d\tau.
\end{equation}

Then

\begin{equation}
\boxed{
p_{\mathrm{eq}}(t)
=
Kp_{\mathrm{ch}}(t)
-
Cm(t)
}
\end{equation}

The term \(m(t)\) represents an exponentially weighted history
of the channel pulse response.

It can be written as

\begin{equation}
m(t)
=
p_{\mathrm{ch}}(t)
*
\left[
e^{-\omega_p t}
u(t)
\right].
\end{equation}

The corresponding transfer function is

\begin{equation}
H_{\mathrm{LP}}(s)
=
\frac{1}
{s+\omega_p}.
\end{equation}

Therefore,

\begin{equation}
\boxed{
m(t)
=
\text{a first-order low-pass filtered version of }
p_{\mathrm{ch}}(t)
}
\end{equation}

apart from a constant gain factor.

% =========================================================
% Differential Equation
% =========================================================
\subsection{Low-Pass Differential Equation}

Starting from

\begin{equation}
m(t)
=
\int_{-\infty}^{t}
p_{\mathrm{ch}}(\tau)
e^{-\omega_p(t-\tau)}
d\tau,
\end{equation}

differentiate with respect to time.

Using Leibniz's rule,

\begin{align}
\frac{dm(t)}{dt}
&=
p_{\mathrm{ch}}(t)
+
\int_{-\infty}^{t}
p_{\mathrm{ch}}(\tau)
\frac{\partial}{\partial t}
e^{-\omega_p(t-\tau)}
d\tau.
\end{align}

Since

\begin{equation}
\frac{\partial}{\partial t}
e^{-\omega_p(t-\tau)}
=
-\omega_p
e^{-\omega_p(t-\tau)},
\end{equation}

we obtain

\begin{align}
\frac{dm(t)}{dt}
&=
p_{\mathrm{ch}}(t)
-
\omega_p
\int_{-\infty}^{t}
p_{\mathrm{ch}}(\tau)
e^{-\omega_p(t-\tau)}
d\tau.
\end{align}

Therefore,

\begin{equation}
\boxed{
\frac{dm(t)}{dt}
+
\omega_p m(t)
=
p_{\mathrm{ch}}(t)
}
\end{equation}

which is the standard first-order low-pass differential equation.

Taking the Laplace transform,

\begin{equation}
sM(s)
+
\omega_pM(s)
=
P_{\mathrm{ch}}(s).
\end{equation}

Thus,

\begin{equation}
\boxed{
\frac{M(s)}
{P_{\mathrm{ch}}(s)}
=
\frac{1}
{s+\omega_p}
}
\end{equation}

which confirms the low-pass behavior.

% =========================================================
% Why Tail Is Reduced
% =========================================================
\subsection{Why the Main Cursor Is Preserved and the Tail Is Reduced}

The equalized response is

\begin{equation}
p_{\mathrm{eq}}(t)
=
Kp_{\mathrm{ch}}(t)
-
Cm(t).
\end{equation}

Consider a rapidly changing portion of the channel pulse response.

During a fast transition,

\begin{equation}
p_{\mathrm{ch}}(t)
\end{equation}

can change rapidly, whereas the low-pass quantity

\begin{equation}
m(t)
\end{equation}

cannot respond instantaneously.

Therefore,

\begin{equation}
m(t)
\end{equation}

is relatively small compared with the instantaneous term, giving

\begin{equation}
p_{\mathrm{eq}}(t)
\approx
Kp_{\mathrm{ch}}(t).
\end{equation}

Thus, the rapidly changing main cursor is preserved or emphasized.

For a slowly varying component, however, \(m(t)\) has enough
time to follow the waveform.

If

\begin{equation}
p_{\mathrm{ch}}(t)=P
\end{equation}

for a sufficiently long time, then

\begin{align}
m(t)
&=
P
\int_0^\infty
e^{-\omega_pu}
du
\\
&=
\frac{P}{\omega_p}.
\end{align}

The output becomes

\begin{equation}
p_{\mathrm{eq}}
=
KP
-
C\frac{P}{\omega_p}.
\end{equation}

Therefore, slowly varying components experience stronger cancellation.

The CTLE can therefore be interpreted as

\begin{equation}
\boxed{
\text{CTLE output}
=
\text{instantaneous signal}
-
\text{low-pass version of the signal}.
}
\end{equation}

Equivalently,

\begin{equation}
\boxed{
\text{high-frequency emphasis}
=
\text{signal}
-
\text{slow component}.
}
\end{equation}

The CTLE does not explicitly identify the main cursor or
post-cursor.

Instead,

\begin{equation}
\boxed{
\text{main cursor}
\rightarrow
\text{more fast-changing content},
}
\end{equation}

while

\begin{equation}
\boxed{
\text{long tail}
\rightarrow
\text{more slowly varying content}.
}
\end{equation}

% =========================================================
% Over Equalization
% =========================================================
\subsection{Over-Equalization}

The CTLE can also overcompensate the channel.

From

\begin{equation}
p_{\mathrm{eq}}(t)
=
Kp_{\mathrm{ch}}(t)
-
Cm(t),
\end{equation}

if

\begin{equation}
Cm(t)
>
Kp_{\mathrm{ch}}(t),
\end{equation}

then

\begin{equation}
\boxed{
p_{\mathrm{eq}}(t)<0.
}
\end{equation}

Thus, an originally positive post-cursor can become negative.

For example, suppose

\begin{equation}
Kp_{\mathrm{ch}}(1\mathrm{UI})
=
0.28,
\end{equation}

and

\begin{equation}
Cm(1\mathrm{UI})
=
0.35.
\end{equation}

Then

\begin{equation}
p_{\mathrm{eq}}(1\mathrm{UI})
=
0.28-0.35
=
-0.07.
\end{equation}

Hence,

\begin{equation}
\boxed{
\text{excessive CTLE boost}
\rightarrow
\text{over-equalization}
\rightarrow
\text{opposite-sign post-cursor}.
}
\end{equation}

Therefore, maximizing CTLE peaking is not the design objective.

Instead, the objective is to maximize eye opening and minimize
the total sampled ISI.

% =========================================================
% Sampling
% =========================================================
\subsection{From Continuous-Time Pulse Response to Discrete Cursors}

After equalization,

\begin{equation}
p_{\mathrm{eq}}(t)
=
p_{\mathrm{ch}}(t)
*
h_{\mathrm{CTLE}}(t).
\end{equation}

Let the symbol period be

\begin{equation}
T_{\mathrm{UI}}.
\end{equation}

Choose a sampling phase

\begin{equation}
t_s.
\end{equation}

The sampled coefficients are

\begin{equation}
\boxed{
c[k]
=
p_{\mathrm{eq}}
\left(
t_s+kT_{\mathrm{UI}}
\right).
}
\end{equation}

The main cursor is

\begin{equation}
\boxed{
c[0]
=
p_{\mathrm{eq}}(t_s).
}
\end{equation}

The first post-cursor is

\begin{equation}
\boxed{
c[1]
=
p_{\mathrm{eq}}
\left(
t_s+T_{\mathrm{UI}}
\right).
}
\end{equation}

The second post-cursor is

\begin{equation}
\boxed{
c[2]
=
p_{\mathrm{eq}}
\left(
t_s+2T_{\mathrm{UI}}
\right).
}
\end{equation}

Similarly, the first pre-cursor is

\begin{equation}
\boxed{
c[-1]
=
p_{\mathrm{eq}}
\left(
t_s-T_{\mathrm{UI}}
\right).
}
\end{equation}

Therefore,

\begin{equation}
\boxed{
c[k]
=
\left[
p_{\mathrm{ch}}
*
h_{\mathrm{CTLE}}
\right]
\left(
t_s+kT_{\mathrm{UI}}
\right).
}
\end{equation}

These are the discrete-time coefficients seen by the sampler
and the DFE.

% =========================================================
% Discrete Channel Derivation
% =========================================================
\subsection{Why Sampled Pulse Response Becomes a Discrete Channel}

Let the transmitted symbol sequence be

\begin{equation}
a[n].
\end{equation}

The transmitted waveform is

\begin{equation}
x(t)
=
\sum_{n=-\infty}^{+\infty}
a[n]
x_{\mathrm{pulse}}
\left(
t-nT_{\mathrm{UI}}
\right).
\end{equation}

After the channel and CTLE,

\begin{equation}
y(t)
=
\sum_{n=-\infty}^{+\infty}
a[n]
p_{\mathrm{eq}}
\left(
t-nT_{\mathrm{UI}}
\right).
\end{equation}

Sample at

\begin{equation}
t
=
mT_{\mathrm{UI}}
+
t_s.
\end{equation}

Then

\begin{align}
y[m]
&=
\sum_n
a[n]
p_{\mathrm{eq}}
\left(
t_s+(m-n)T_{\mathrm{UI}}
\right).
\end{align}

Using

\begin{equation}
c[k]
=
p_{\mathrm{eq}}
\left(
t_s+kT_{\mathrm{UI}}
\right),
\end{equation}

we obtain

\begin{equation}
\boxed{
y[m]
=
\sum_n
a[n]c[m-n].
}
\end{equation}

Equivalently,

\begin{equation}
y[m]
=
a[m]c[0]
+
a[m-1]c[1]
+
a[m-2]c[2]
+\cdots.
\end{equation}

Thus,

\begin{equation}
\boxed{
c[0]
=
\text{main cursor},
}
\end{equation}

and

\begin{equation}
\boxed{
c[1],c[2],\ldots
=
\text{post-cursor ISI coefficients}.
}
\end{equation}

% =========================================================
% Sampling and Convolution
% =========================================================
\subsection{Sampling Before and After Convolution}

For two continuous-time signals

\begin{equation}
x(t)
\end{equation}

and

\begin{equation}
h(t),
\end{equation}

the continuous-time convolution is

\begin{equation}
y(t)
=
(x*h)(t)
=
\int_{-\infty}^{+\infty}
x(\tau)
h(t-\tau)
d\tau.
\end{equation}

Sampling after convolution gives

\begin{equation}
\boxed{
y[n]
=
y(nT_s)
=
\int_{-\infty}^{+\infty}
x(\tau)
h(nT_s-\tau)
d\tau.
}
\end{equation}

If the signals are sampled first,

\begin{equation}
x[n]
=
x(nT_s),
\end{equation}

and

\begin{equation}
h[n]
=
h(nT_s),
\end{equation}

then the discrete convolution is

\begin{equation}
\boxed{
y_d[n]
=
\sum_{k=-\infty}^{+\infty}
x(kT_s)
h((n-k)T_s).
}
\end{equation}

In general,

\begin{equation}
\boxed{
\operatorname{sample}
\left\{
x*h
\right\}
\neq
\operatorname{sample}\{x\}
*
\operatorname{sample}\{h\}.
}
\end{equation}

The reason is that continuous-time convolution uses every
value of the waveform between sampling instants.

The exact sampled continuous-time convolution can be written as

\begin{equation}
y(nT_s)
=
\sum_k
\int_{kT_s}^{(k+1)T_s}
x(\tau)
h(nT_s-\tau)
d\tau.
\end{equation}

If \(T_s\) is sufficiently small, each integral can be approximated by

\begin{equation}
\int_{kT_s}^{(k+1)T_s}
x(\tau)
h(nT_s-\tau)
d\tau
\approx
T_s
x(kT_s)
h((n-k)T_s).
\end{equation}

Therefore,

\begin{equation}
\boxed{
y(nT_s)
\approx
T_s
\sum_k
x(kT_s)
h((n-k)T_s).
}
\end{equation}

Thus, high-rate discrete convolution can approximate
continuous-time convolution.

However, sampling only once per UI before convolution generally
loses important intra-UI waveform information.

% =========================================================
% Numerical Procedure
% =========================================================
\subsection{Recommended Numerical Procedure}

Given a measured or simulated channel pulse response
\(p_{\mathrm{ch}}(t)\), a practical numerical flow is:

\begin{enumerate}[label=\arabic*.]

    \item Obtain the CTLE impulse response

    \begin{equation}
    h_{\mathrm{CTLE}}(t).
    \end{equation}

    \item Choose a fine time step

    \begin{equation}
    \Delta t
    \ll
    T_{\mathrm{UI}}.
    \end{equation}

    \item Sample the channel pulse response

    \begin{equation}
    p_{\mathrm{ch}}[n]
    =
    p_{\mathrm{ch}}(n\Delta t).
    \end{equation}

    \item Sample the CTLE impulse response

    \begin{equation}
    h_{\mathrm{CTLE}}[n]
    =
    h_{\mathrm{CTLE}}(n\Delta t).
    \end{equation}

    \item Perform numerical convolution

    \begin{equation}
    p_{\mathrm{eq}}[n]
    \approx
    \Delta t
    \left(
    p_{\mathrm{ch}}
    *
    h_{\mathrm{CTLE}}
    \right)[n].
    \end{equation}

    \item Select the sampling phase \(t_s\).

    \item Sample once every UI

    \begin{equation}
    c[k]
    =
    p_{\mathrm{eq}}
    \left(
    t_s+kT_{\mathrm{UI}}
    \right).
    \end{equation}

    \item Extract

    \begin{equation}
    c[0],
    c[1],
    c[2],
    \ldots
    \end{equation}

    as the main cursor and post-cursor coefficients.

\end{enumerate}

% =========================================================
% Summary
% =========================================================
\section{Summary}

The CTLE equalization process can be summarized as

\begin{equation}
\boxed{
p_{\mathrm{ch}}(t)
\xrightarrow{*\,h_{\mathrm{CTLE}}(t)}
p_{\mathrm{eq}}(t)
\xrightarrow{\text{sample every UI}}
c[k].
}
\end{equation}

The CTLE impulse response can be written as

\begin{equation}
\boxed{
h_{\mathrm{CTLE}}(t)
=
K\delta(t)
-
Ce^{-\omega_p t}u(t).
}
\end{equation}

The equalized pulse response is

\begin{equation}
\boxed{
p_{\mathrm{eq}}(t)
=
Kp_{\mathrm{ch}}(t)
-
C
\int_{-\infty}^{t}
p_{\mathrm{ch}}(\tau)
e^{-\omega_p(t-\tau)}
d\tau.
}
\end{equation}

The first term represents the instantaneous component, while
the second term subtracts an exponentially weighted low-pass
version of the waveform.

Therefore,

\begin{equation}
\boxed{
\text{fast-changing waveform components}
\rightarrow
\text{relatively emphasized},
}
\end{equation}

while

\begin{equation}
\boxed{
\text{slowly varying channel tails}
\rightarrow
\text{relatively suppressed}.
}
\end{equation}

Excessive peaking may cause

\begin{equation}
\boxed{
\text{over-equalization}
\rightarrow
\text{negative post-cursors}.
}
\end{equation}

The sampled equalized pulse response

\begin{equation}
\boxed{
c[k]
=
p_{\mathrm{eq}}
\left(
t_s+kT_{\mathrm{UI}}
\right)
}
\end{equation}

provides the discrete-time cursor coefficients used for
ISI and DFE analysis.

% =========================================================
% Placeholder Chapters
% =========================================================
\chapter{Decision-Feedback Equalization}

This chapter will discuss DFE architecture, feedback timing,
critical-path constraints, and PAM-3 DFE implementation.

\chapter{PAM-3 Receiver Architecture}

This chapter will discuss PAM-3 slicers, upper and lower
thresholds, VREFH/VREFL, receiver offset calibration, and
tri-level decision logic.

\chapter{Clocking and Timing}

This chapter will discuss sampling clocks, forwarded clocks,
clock phase control, duty-cycle correction, and timing margin.

\chapter{GDDR7 Receiver Design}

This chapter will discuss GDDR7 receiver architecture,
training, VREF calibration, CTLE, DFE, slicers, and timing.

% =========================================================
% Back Matter
% =========================================================
\backmatter

\chapter*{Notes}

Additional design notes, references, and derivations can be
added here.

\end{document}